\chapter{Algebra}

\begin{formula}[Basic Factorization Identities]
For elements of a commutative ring,
\[
 a^2-b^2=(a-b)(a+b), \qquad
 a^3\pm b^3=(a\pm b)(a^2\mp ab+b^2),
\]
where the two signs in the second identity are paired consistently. More generally, for $n\ge1$,
\[
 a^n-b^n=(a-b)\sum_{k=0}^{n-1}a^{n-1-k}b^k.
\]
If $n$ is odd, then
\[
 a^n+b^n=(a+b)\sum_{k=0}^{n-1}(-1)^k a^{n-1-k}b^k.
\]
\end{formula}

\begin{formula}[Three-Variable Cubic Identity]
For $a, b, c$ in a commutative ring,
\[
 a^3+b^3+c^3-3abc
 =(a+b+c)(a^2+b^2+c^2-ab-bc-ca).
\]
In particular, $a+b+c=0$ implies $a^3+b^3+c^3=3abc$.
\end{formula}

\begin{formula}[Square and Fourth-Power Sums]
For real $a, b, c$,
\[
 a^2+b^2+c^2-ab-bc-ca
 =\frac12\bigl((a-b)^2+(b-c)^2+(c-a)^2\bigr),
\]
\[
 a^4+b^4+c^4-\sum_{\mathrm{cyc}}a^2b^2
 =\frac12\bigl((a^2-b^2)^2+(b^2-c^2)^2+(c^2-a^2)^2\bigr).
\]
Here $\sum_{\mathrm{cyc}}a^2b^2=a^2b^2+b^2c^2+c^2a^2$.
\end{formula}

\begin{definition}[Symmetric and Homogeneous Polynomials]
A polynomial $P(x_1, \dots, x_n)$ is \emph{symmetric} if it is unchanged by every permutation of its variables. It is \emph{homogeneous of degree $d$} if every nonzero monomial has total degree $d$, equivalently
\[
 P(tx_1, \dots, tx_n)=t^d P(x_1, \dots, x_n)
\]
for every scalar $t$ in an infinite coefficient field.
\end{definition}

\begin{formula}[Elementary Symmetric Polynomials]
For $x_1, \dots, x_n$ and $0\le k\le n$, put $e_0=1$ and
\[
 e_k=\sum_{1\le i_1<\cdots<i_k\le n}x_{i_1}\cdots x_{i_k}.
\]
Then
\[
 \prod_{i=1}^n(t-x_i)=\sum_{k=0}^n(-1)^k e_k t^{n-k}.
\]
\end{formula}

\begin{theorem}[Fundamental Theorem of Symmetric Polynomials]
Every symmetric polynomial in $x_1, \dots, x_n$ over a commutative ring $R$ can be expressed uniquely as a polynomial in $e_1, \dots, e_n$ with coefficients in $R$.
\end{theorem}

\begin{formula}[Quadratic Equation]
For $a, b, c\in\C$ with $a\ne0$, the solutions of $ax^2+bx+c=0$ are
\[
 x=\frac{-b\pm\sqrt{b^2-4ac}}{2a}.
\]
The discriminant is $\Delta=b^2-4ac$. For real coefficients, its sign determines whether there are two distinct, one repeated, or no real roots.
\end{formula}

\begin{formula}[Two-Variable Symmetric Expressions]
If $s=x+y$ and $p=xy$, then $x, y$ are roots of $t^2-st+p$, and
\[
 x^2+y^2=s^2-2p, \quad
 x^3+y^3=s^3-3ps, \quad
 (x-y)^2=s^2-4p.
\]
\end{formula}


\begin{theorem}[Vieta's Relations for Two and Three Roots]
If $x, y$ are roots of $t^2-st+p$, then $x+y=s$ and $xy=p$. If $x, y, z$ are roots of $t^3-at^2+bt-c$, then
\[
 x+y+z=a, \quad xy+yz+zx=b, \quad xyz=c.
\]
\end{theorem}

\begin{formula}[Sums of Powers of Integers]
For $n\in\N$,
\[
 \sum_{k=1}^n k=\frac{n(n+1)}2, \qquad
 \sum_{k=1}^n k^2=\frac{n(n+1)(2n+1)}6,
\]
\[
 \sum_{k=1}^n k^3=\left(\frac{n(n+1)}2\right)^2.
\]
\end{formula}

\begin{definition}[Polynomial, Degree, and Multiplicity]
For a field $F$, a polynomial $f(x)=a_0+\cdots+a_nx^n\in F[x]$ with $a_n\ne0$ has degree $n$. A root $r\in F$ has multiplicity $m\ge1$ if $(x-r)^m\mid f(x)$ but $(x-r)^{m+1}\nmid f(x)$.
\end{definition}

\begin{theorem}[Polynomial Division and Remainder Theorems]
For $f, g\in F[x]$ with $g\ne0$, there are unique $q, r\in F[x]$ satisfying
\[
 f=qg+r, \qquad r=0\ \text{or}\ \deg r<\deg g.
\]
In particular, the remainder on division of $f(x)$ by $x-a$ is $f(a)$, and
\[
 f(a)=0\quad\Longleftrightarrow\quad (x-a)\mid f(x).
\]
\end{theorem}


\begin{theorem}[Polynomial B\'ezout Identity]
Let $F$ be a field and let $f, g\in F[x]$, not both zero. There are polynomials $u, v\in F[x]$ such that
\[
 uf+vg=\gcd(f, g),
\]
where the greatest common divisor is chosen monic. In particular, $f$ and $g$ are coprime if and only if $uf+vg=1$ for some $u, v\in F[x]$.
\end{theorem}



\begin{proposition}[Root-Counting Principle]
A nonzero polynomial of degree $d$ over a field has at most $d$ distinct roots. Consequently, two polynomials of degree at most $d$ that agree at $d+1$ distinct points are identical.
\end{proposition}

\begin{theorem}[Vi\`ete's Formulas]
Let $f(x)=a_nx^n+\cdots+a_0=a_n\prod_{i=1}^n(x-r_i)$ over $\C$, with $a_n\ne0$, where the roots are counted with multiplicity. Then
\[
 e_k(r_1, \dots, r_n)=(-1)^k\frac{a_{n-k}}{a_n}
 \qquad (1\le k\le n).
\]
In particular, $\sum_i r_i=-a_{n-1}/a_n$ and $\prod_i r_i=(-1)^na_0/a_n$.
\end{theorem}

\begin{formula}[Newton's Identities]
Let $p_k=\sum_{i=1}^n r_i^k$ and let $e_k$ be the elementary symmetric polynomials in $r_1, \dots, r_n$. For $1\le k\le n$,
\[
 p_k-e_1p_{k-1}+e_2p_{k-2}-\cdots
 +(-1)^{k-1}e_{k-1}p_1+(-1)^k k e_k=0.
\]
For $k>n$, the same identity holds with $e_j=0$ for $j>n$.
\end{formula}

\begin{theorem}[Lagrange Interpolation]
Let $x_0, \dots, x_n$ be distinct elements of a field $F$. There is a unique polynomial $P\in F[x]$ of degree at most $n$ with $P(x_i)=y_i$ for all $i$, namely
\[
 P(x)=\sum_{i=0}^n y_i\prod_{\substack{0\le j\le n\\j\ne i}}
 \frac{x-x_j}{x_i-x_j}.
\]
\end{theorem}

\begin{theorem}[Rational Root Theorem]
Let $f(x)=a_nx^n+\cdots+a_0\in\Z[x]$ with $a_n\ne0$. If $u/v$ is a root with $u, v\in\Z$, $v>0$, and $\gcd(u, v)=1$, then $u\mid a_0$ and $v\mid a_n$.
\end{theorem}

\begin{theorem}[Conjugate Root Principles]
For a polynomial with real coefficients, every nonreal complex root occurs together with its complex conjugate, with the same multiplicity. For a polynomial with rational coefficients, a root of the form $a+b\sqrt d$ with $a, b\in\Q$, $b\ne0$, and $d\in\Q$ not a square in $\Q$ occurs together with $a-b\sqrt d$.
\end{theorem}

\begin{theorem}[Eisenstein's Irreducibility Criterion]
Let $f(x)=a_nx^n+\cdots+a_0\in\Z[x]$ be primitive. If a prime $p$ satisfies
\[
 p\mid a_0, \dots, a_{n-1}, \qquad p\nmid a_n,
 \qquad p^2\nmid a_0,
\]
then $f$ is irreducible over $\Q$. The criterion can also be applied after a substitution $x\mapsto x+k$ with $k\in\Z$.
\end{theorem}

\begin{theorem}[Descartes' Rule of Signs]
Let $f(x)\in\R[x]$ be a nonzero polynomial. The number of its positive real roots, counted with multiplicity, is at most the number of sign changes in the sequence of its nonzero coefficients and differs from that number by an even integer. Applying the same rule to $f(-x)$ bounds its negative real roots.
\end{theorem}

\begin{theorem}[Cauchy's Root Bound]
If $f(z)=a_nz^n+\cdots+a_0$ with $a_n\ne0$, then every complex root $z$ satisfies
\[
 |z|\le 1+\max_{0\le k<n}\left|\frac{a_k}{a_n}\right|.
\]
\end{theorem}

\begin{definition}[Irreducible Polynomial and Minimal Polynomial]
A nonconstant $f\in F[x]$ is \emph{irreducible over $F$} if it cannot be written as a product of two polynomials of positive degree in $F[x]$. If $\alpha$ is algebraic over $F$, its \emph{minimal polynomial} $m_{\alpha, F}\in F[x]$ is the unique monic irreducible polynomial over $F$ satisfying $m_{\alpha, F}(\alpha)=0$.
\end{definition}

\begin{theorem}[Fundamental Theorem of Algebra]
Every nonconstant polynomial in $\C[x]$ has a root in $\C$. Equivalently, every polynomial of degree $n\ge1$ over $\C$ factors into $n$ linear factors, counted with multiplicity.
\end{theorem}

\begin{theorem}[Combinatorial Nullstellensatz, Coefficient Form]
Let $F$ be a field and let $P\in F[x_1, \dots, x_n]$ have total degree $t_1+\cdots+t_n$, where $t_i\in\Nzero$. If the coefficient of $x_1^{t_1}\cdots x_n^{t_n}$ is nonzero, then for any sets $S_i\subseteq F$ with $|S_i|>t_i$, there exist $s_i\in S_i$ such that $P(s_1, \dots, s_n)\ne0$.
\end{theorem}


\begin{theorem}[Gauss's Lemma for Polynomials]
If $f, g\in\Z[x]$ are primitive polynomials, then $fg$ is primitive. In particular, a primitive polynomial in $\Z[x]$ is irreducible over $\Q$ if and only if it is irreducible over $\Z$.
\end{theorem}

\begin{theorem}[Irreducibility by Reduction Modulo a Prime]
Let $f\in\Z[x]$ have positive degree and leading coefficient not divisible by a prime $p$. If its reduction $\overline f\in\mathbb F_p[x]$ has the same degree and is irreducible over $\mathbb F_p$, then $f$ is irreducible over $\Q$ provided $f$ is primitive.
\end{theorem}

\begin{proposition}[Derivative and Multiplicity of a Root]
Let $f\in F[x]$, where $F$ has characteristic $0$. An element $a\in F$ is a root of multiplicity exactly $m$ if and only if
\[
 f(a)=f'(a)=\cdots=f^{(m-1)}(a)=0,
 \qquad f^{(m)}(a)\ne0.
\]
\end{proposition}

\begin{theorem}[Arithmetic--Geometric Mean Inequality]
For $x_1, \dots, x_n\ge0$,
\[
 \frac{x_1+\cdots+x_n}{n}\ge (x_1\cdots x_n)^{1/n}.
\]
Equality holds exactly when $x_1=\cdots=x_n$. More generally, for positive weights $w_i$ with $\sum_iw_i=1$,
\[
 \sum_{i=1}^n w_ix_i\ge\prod_{i=1}^n x_i^{w_i},
\]
with equality when all $x_i>0$ are equal.
\end{theorem}

\begin{theorem}[Cauchy--Schwarz Inequality]
For real $a_i, b_i$,
\[
 \left(\sum_{i=1}^n a_ib_i\right)^2
 \le \left(\sum_{i=1}^n a_i^2\right)
       \left(\sum_{i=1}^n b_i^2\right).
\]
Equality holds exactly when the two vectors are linearly dependent. For $x_i>0$,
\[
 \sum_{i=1}^n\frac{a_i^2}{x_i}
 \ge\frac{(\sum_{i=1}^n a_i)^2}{\sum_{i=1}^n x_i}.
\]
\end{theorem}


\begin{theorem}[Cauchy--Schwarz Inequality, Engel Form]
For $x_1, \dots, x_n\in\R$ and $a_1, \dots, a_n>0$,
\[
 \sum_{i=1}^n\frac{x_i^2}{a_i}\ge
 \frac{(x_1+\cdots+x_n)^2}{a_1+\cdots+a_n}.
\]
Equality holds if and only if the ratios $x_i/a_i$ are all equal.
\end{theorem}

\begin{theorem}[Tangent-Line Inequality for Convex Functions]
If $f$ is differentiable and convex on an open interval $I$, then for all $x, y\in I$,
\[
 f(y)\ge f(x)+f'(x)(y-x).
\]
If $f$ is strictly convex, equality holds only for $x=y$.
\end{theorem}

\begin{theorem}[Harmonic--Geometric--Arithmetic--Quadratic Means]
For positive $x_1, \dots, x_n$,
\[
 \frac{n}{\sum_i1/x_i}
 \le\left(\prod_i x_i\right)^{1/n}
 \le\frac{\sum_i x_i}{n}
 \le\sqrt{\frac{\sum_i x_i^2}{n}}.
\]
Equality throughout holds exactly when all $x_i$ are equal.
\end{theorem}

\begin{theorem}[H\"older's Inequality]
Let $p, q>1$ satisfy $1/p+1/q=1$. For nonnegative $a_i, b_i$,
\[
 \sum_{i=1}^n a_ib_i\le
 \left(\sum_{i=1}^n a_i^p\right)^{1/p}
 \left(\sum_{i=1}^n b_i^q\right)^{1/q}.
\]
For nonzero sequences, equality holds exactly when $(a_i^p)$ and $(b_i^q)$ are proportional.
\end{theorem}

\begin{theorem}[Minkowski's Inequality]
For $p\ge1$ and real (or complex) sequences $(a_i), (b_i)$,
\[
 \left(\sum_i|a_i+b_i|^p\right)^{1/p}
 \le\left(\sum_i|a_i|^p\right)^{1/p}
 +\left(\sum_i|b_i|^p\right)^{1/p}.
\]
\end{theorem}

\begin{theorem}[Jensen's Inequality]
Let $f$ be convex on an interval $I$, and let $x_i\in I$, $w_i\ge0$, $\sum_iw_i=1$. Then
\[
 f\!\left(\sum_iw_ix_i\right)\le\sum_iw_if(x_i).
\]
For a strictly convex $f$ and positive weights, equality holds exactly when all $x_i$ are equal.
\end{theorem}

\begin{theorem}[Bernoulli's Inequality]
For $x\ge-1$ and real $r\ge1$,
\[
 (1+x)^r\ge1+rx.
\]
For $0\le r\le1$ the inequality reverses when $x>-1$.
\end{theorem}

\begin{theorem}[Rearrangement Inequality]
If $a_1\le\cdots\le a_n$ and $b_1\le\cdots\le b_n$ are real, then for any permutation $\sigma\in S_n$,
\[
 \sum_{i=1}^n a_ib_{n+1-i}
 \le\sum_{i=1}^n a_ib_{\sigma(i)}
 \le\sum_{i=1}^n a_ib_i.
\]
\end{theorem}

\begin{theorem}[Chebyshev's Sum Inequality]
If $(a_i)$ and $(b_i)$ are similarly sorted real sequences, then
\[
 \frac1n\sum_{i=1}^n a_ib_i
 \ge\left(\frac1n\sum_{i=1}^n a_i\right)
       \left(\frac1n\sum_{i=1}^n b_i\right).
\]
The inequality reverses if the two sequences are oppositely sorted.
\end{theorem}

\begin{theorem}[Schur's Inequality]
For $a, b, c\ge0$ and $r>0$ (or for $a, b, c>0$ and $r=0$),
\[
 \sum_{\mathrm{cyc}}a^r(a-b)(a-c)\ge0.
\]
For $r=1$, this becomes
\[
 a(a-b)(a-c)+b(b-c)(b-a)+c(c-a)(c-b)\ge0.
\]
\end{theorem}

\begin{theorem}[Nesbitt's Inequality]
For $a, b, c>0$,
\[
 \frac{a}{b+c}+\frac{b}{c+a}+\frac{c}{a+b}\ge\frac32.
\]
Equality holds exactly when $a=b=c$.
\end{theorem}

\begin{theorem}[Young's Inequality for Products]
For $a, b\ge0$ and $p, q>1$ with $1/p+1/q=1$,
\[
 ab\le\frac{a^p}{p}+\frac{b^q}{q}.
\]
Equality holds exactly when $a^p=b^q$.
\end{theorem}

\begin{definition}[Majorization]
For vectors $x, y\in\R^n$, let $x_{[1]}\ge\cdots\ge x_{[n]}$ and $y_{[1]}\ge\cdots\ge y_{[n]}$ be their coordinates in decreasing order. We say that $x$ \emph{majorizes} $y$ if
\[
 \sum_{i=1}^k x_{[i]}\ge\sum_{i=1}^k y_{[i]}
 \quad(1\le k<n), \qquad \sum_i x_i=\sum_i y_i.
\]
\end{definition}

\begin{theorem}[Muirhead's Inequality]
Let $a_1, \dots, a_n\ge0$. For a nonnegative integer $n$-tuple $\alpha$, let $m_\alpha(a)$ denote the average of $a_1^{\alpha_{\sigma(1)}}\cdots a_n^{\alpha_{\sigma(n)}}$ over all $\sigma\in S_n$. If $\alpha$ majorizes $\beta$, then
\[
 m_\alpha(a)\ge m_\beta(a).
\]
\end{theorem}


\begin{theorem}[Power Mean Inequality]
For positive numbers $x_1, \dots, x_n$ and real $r$, let
\[
 M_r(x)=\begin{cases}
 (\frac1n\sum_i x_i^r)^{1/r}, &r\ne0, \\
 (\prod_i x_i)^{1/n}, &r=0.
 \end{cases}
\]
If $r<s$, then $M_r(x)\le M_s(x)$, with equality exactly when all $x_i$ are equal.
\end{theorem}

\begin{theorem}[Maclaurin's Inequalities]
For positive $x_1, \dots, x_n$, let $e_k$ be the elementary symmetric sums and put
\[
 E_k=\frac{e_k}{\binom nk}\qquad(1\le k\le n).
\]
Then
\[
 E_1\ge E_2^{1/2}\ge\cdots\ge E_n^{1/n},
\]
with equality throughout exactly when $x_1=\cdots=x_n$.
\end{theorem}

\begin{theorem}[Karamata's Inequality]
Let $x, y\in I^n$ for an interval $I$, and suppose $x$ majorizes $y$. For a convex function $f\colon I\to\R$,
\[
 \sum_{i=1}^n f(x_i)\ge\sum_{i=1}^n f(y_i).
\]
If $f$ is strictly convex, equality holds exactly when $x$ is a permutation of $y$.
\end{theorem}

\begin{definition}[Sequence and Linear Recurrence]
A sequence $(a_n)$ over a field $F$ is a function from an index set such as $\Nzero$ to $F$. It satisfies a linear recurrence of order $k$ with constant coefficients if
\[
 a_{n+k}=c_1a_{n+k-1}+\cdots+c_ka_n
\]
for fixed $c_1, \dots, c_k\in F$, with $c_k\ne0$.
\end{definition}

\begin{formula}[Arithmetic and Geometric Progressions]
For $n\ge1$, an arithmetic progression $a_n=a_1+(n-1)d$ satisfies
\[
 \sum_{k=1}^n a_k=\frac{n(a_1+a_n)}2.
\]
A geometric progression $a_n=a_1r^{n-1}$ satisfies
\[
 \sum_{k=1}^n a_1r^{k-1}=
 \begin{cases}a_1(1-r^n)/(1-r), &r\ne1, \\na_1, &r=1.\end{cases}
\]
\end{formula}

\begin{theorem}[Linear Recurrences and Characteristic Roots]
Suppose a sequence over $\C$ satisfies
\[
 a_{n+k}=c_1a_{n+k-1}+\cdots+c_ka_n
\]
with $c_k\ne0$. If the characteristic polynomial
\[
 P(t)=t^k-c_1t^{k-1}-\cdots-c_k
\]
has distinct roots $r_i$ of multiplicities $m_i$, then every solution is of the form
\[
 a_n=\sum_i Q_i(n)r_i^n, \qquad \deg Q_i<m_i.
\]
The constants in the polynomials $Q_i$ are determined uniquely by $k$ consecutive initial values.
\end{theorem}

\begin{formula}[Fibonacci Numbers and Binet's Formula]
Let $F_0=0$, $F_1=1$, and $F_{n+2}=F_{n+1}+F_n$. Then
\[
 F_n=\frac{\varphi^n-\psi^n}{\sqrt5}, \qquad
 \varphi=\frac{1+\sqrt5}2, \quad \psi=\frac{1-\sqrt5}2.
\]
For $m, n\ge1$,
\[
 F_{m+n}=F_mF_{n+1}+F_{m-1}F_n,
 \qquad F_{n+1}F_{n-1}-F_n^2=(-1)^n.
\]
\end{formula}

\begin{formula}[Telescoping Identities]
For any sequence $(a_k)$ and $m\le n$,
\[
 \sum_{k=m}^n(a_{k+1}-a_k)=a_{n+1}-a_m.
\]
If $a_k, a_{k+1}\ne0$,
\[
 \frac{a_{k+1}-a_k}{a_ka_{k+1}}
 =\frac1{a_k}-\frac1{a_{k+1}}.
\]
\end{formula}


\begin{formula}[Finite Differences and Newton's Formula]
For a sequence $(a_n)_{n\ge0}$, define $(\Delta a)_n=a_{n+1}-a_n$. If $\Delta^{d+1}a=0$, then
\[
 a_n=\sum_{k=0}^d \binom nk(\Delta^k a)_0
 \qquad(n\in\Nzero).
\]
In particular, if $a_n=P(n)$ for a polynomial $P$ of degree at most $d$, then $\Delta^{d+1}a=0$.
\end{formula}

\begin{formula}[First-Order Linear Recurrence]
If $a_{n+1}=ra_n+b_n$ for $n\ge0$, with fixed $r\in\C$, then
\[
 a_n=r^na_0+\sum_{j=0}^{n-1}r^{n-1-j}b_j
 \qquad(n\ge1).
\]
If $b_n=b$ is constant and $r\ne1$, this becomes
\[
 a_n=r^na_0+b\frac{r^n-1}{r-1}.
\]
\end{formula}

\begin{definition}[Injective, Surjective, and Bijective Functions]
A function $f\colon X\to Y$ is injective if $f(x)=f(y)$ implies $x=y$, surjective if every $y\in Y$ equals $f(x)$ for some $x\in X$, and bijective if both properties hold. A bijection has an inverse function $f^{-1}\colon Y\to X$.
\end{definition}

\begin{definition}[Additive and Multiplicative Functions]
A function $f\colon\R\to\R$ is additive if $f(x+y)=f(x)+f(y)$ for all $x, y\in\R$. A function $g\colon\R\to\R$ is multiplicative if $g(xy)=g(x)g(y)$ for all $x, y\in\R$. These notions are distinct from multiplicative arithmetic functions on $\N$.
\end{definition}

\begin{theorem}[Cauchy's Additive Functional Equation]
If $f\colon\R\to\R$ satisfies $f(x+y)=f(x)+f(y)$ for all $x, y\in\R$, then $f(q)=qf(1)$ for all $q\in\Q$. If, in addition, $f$ is continuous at one point, monotone on an interval, Lebesgue measurable, or bounded on an interval of positive length, then
\[
 f(x)=cx\qquad(x\in\R)
\]
for some $c\in\R$.
\end{theorem}

\begin{theorem}[Jensen's Functional Equation]
If $f\colon\R\to\R$ satisfies
\[
 f\!\left(\frac{x+y}{2}\right)=\frac{f(x)+f(y)}2
 \qquad(x, y\in\R),
\]
then $f(x)=A(x)+b$ for an additive function $A\colon\R\to\R$ and a constant $b\in\R$. If $f$ is continuous at one point, then $f(x)=cx+b$ for real $b, c$.
\end{theorem}

\begin{theorem}[Continuous Exponential Equation]
If a continuous function $f\colon\R\to(0, \infty)$ satisfies
\[
 f(x+y)=f(x)f(y)\qquad(x, y\in\R),
\]
then $f(x)=e^{cx}$ for some $c\in\R$.
\end{theorem}

\begin{remark}
In a functional equation, the domain, codomain, and any regularity assumptions are part of the statement. The same formal equation may have additional nonstandard solutions without continuity, measurability, or related hypotheses.
\end{remark}


\begin{theorem}[Polynomial Functional Equations]
Suppose $P, Q\in F[x]$ over an infinite field $F$ satisfy $P(t)=Q(t)$ for every $t\in F$. Then $P=Q$ as polynomials, so their corresponding coefficients are equal.
\end{theorem}

\begin{formula}[Polar Form and De Moivre's Formula]
For $z=r(\cos\theta+i\sin\theta)$ with $r\ge0$ and $n\in\Z$ (when $n<0$, assume $r>0$),
\[
 z^n=r^n(\cos(n\theta)+i\sin(n\theta)).
\]
The $n$th roots of a nonzero complex number $re^{i\theta}$ are
\[
 r^{1/n}e^{i(\theta+2\pi k)/n}, \qquad k=0, \dots, n-1.
\]
\end{formula}

\begin{definition}[Roots of Unity]
For $n\ge1$, the $n$th roots of unity are the solutions of $z^n=1$. They form a cyclic group of order $n$, generated by
\[
 \zeta_n=e^{2\pi i/n}.
\]
An $n$th root of unity is \emph{primitive} if its multiplicative order is $n$.
\end{definition}

\begin{formula}[Sums of Roots of Unity]
For $n\ge2$ and any integer $m$,
\[
 \sum_{k=0}^{n-1}\zeta_n^{mk}=
 \begin{cases}n, &n\mid m, \\0, &n\nmid m.\end{cases}
\]
The same identity holds for any primitive $n$th root in place of $\zeta_n$.
\end{formula}

\begin{formula}[Roots-of-Unity Filter]
If $P(x)=\sum_{k=0}^N a_kx^k\in\C[x]$ and $r\in\Z$, then
\[
 \sum_{\substack{0\le k\le N\\k\equiv r\ (\mathrm{mod}\ n)}}a_k
 =\frac1n\sum_{j=0}^{n-1}\zeta_n^{-rj}P(\zeta_n^j).
\]
\end{formula}

\begin{theorem}[Cyclotomic Factorization]
For $n\ge1$, let $\Phi_n(x)$ be the monic polynomial whose roots are the primitive $n$th roots of unity. Then $\Phi_n(x)\in\Z[x]$, $\deg\Phi_n=\varphi(n)$, and
\[
 x^n-1=\prod_{d\mid n}\Phi_d(x).
\]
Each $\Phi_n$ is irreducible over $\Q$. Here $\varphi(n)$ denotes the number of integers $1\le k\le n$ relatively prime to $n$.
\end{theorem}

\begin{definition}[Group and Subgroup]
A group $(G, \ast)$ is a set with an associative binary operation, an identity $e$, and a two-sided inverse for each element. A subset $H\subseteq G$ is a subgroup, written $H\le G$, if it is a group under the restricted operation. The group is \emph{abelian} if $a\ast b=b\ast a$ for all $a, b\in G$.
\end{definition}

\begin{definition}[Cyclic and Permutation Groups]
A group generated by one element $g$ is cyclic, written $\langle g\rangle$. Write $C_n$ for the cyclic group of order $n$ and $S_n$ for the permutation group on $n$ symbols. The dihedral group $D_n$ is the symmetry group of a regular $n$-gon and has order $2n$.
\end{definition}

\begin{theorem}[Lagrange's Theorem]
If $H\le G$ and $G$ is finite, then
\[
 |G|=[G:H]|H|.
\]
Consequently, $|H|\mid|G|$ and the order $\ord(g)$ of any $g\in G$ divides $|G|$; in particular, $g^{|G|}=e$.
\end{theorem}

\begin{theorem}[Cauchy's Theorem for Finite Groups]
If a prime $p$ divides the order of a finite group $G$, then $G$ has an element of order $p$.
\end{theorem}

\begin{theorem}[Structure of Finite Cyclic Groups]
If $G=\langle g\rangle$ has order $n$, then
\[
 \ord(g^k)=\frac{n}{\gcd(n, k)}.
\]
For each positive divisor $d\mid n$ there is exactly one subgroup of $G$ of order $d$, and $G$ has exactly $\varphi(n)$ generators, where $\varphi$ is Euler's totient function.
\end{theorem}

\begin{definition}[Homomorphism, Kernel, and Quotient Group]
A group homomorphism $\varphi\colon G\to H$ satisfies $\varphi(xy)=\varphi(x)\varphi(y)$. Its kernel is $\ker\varphi=\{g\in G\mid\varphi(g)=e_H\}$. A subgroup $N\le G$ is normal if $gNg^{-1}=N$ for all $g\in G$; in that case $G/N$ is the quotient group of cosets.
\end{definition}

\begin{theorem}[First Isomorphism Theorem for Groups]
For any group homomorphism $\varphi\colon G\to H$,
\[
 G/\ker\varphi\cong\im\varphi.
\]
\end{theorem}

\begin{definition}[Group Action, Orbit, and Stabilizer]
A left action of $G$ on $X$ is a map $(g, x)\mapsto g\cdot x$ satisfying $e\cdot x=x$ and $(gh)\cdot x=g\cdot(h\cdot x)$. The orbit of $x$ is $G\cdot x=\{g\cdot x\mid g\in G\}$, and the stabilizer is $\Stab_G(x)=\{g\in G\mid g\cdot x=x\}$.
\end{definition}

\begin{theorem}[Orbit--Stabilizer Theorem]
If a finite group $G$ acts on a set $X$, then for each $x\in X$,
\[
 |G\cdot x|=[G:\Stab_G(x)]
 =\frac{|G|}{|\Stab_G(x)|}.
\]
\end{theorem}

\begin{theorem}[Burnside's Lemma; Cauchy--Frobenius Formula]
If a finite group $G$ acts on a finite set $X$, the number of orbits is
\[
 |G\backslash X|=\frac1{|G|}\sum_{g\in G}|\Fix(g)|,
\]
where $\Fix(g)=\{x\in X\mid g\cdot x=x\}$.
\end{theorem}

\begin{theorem}[Class Equation]
For a finite group $G$, choose representatives $x_1, \dots, x_r$ of the conjugacy classes not contained in its center $Z(G)$. Then
\[
 |G|=|Z(G)|+\sum_{i=1}^r[G:C_G(x_i)].
\]
\end{theorem}

\begin{theorem}[Sylow Theorems]
Let $|G|=p^km$ where $p$ is prime and $p\nmid m$. Then $G$ has a subgroup of order $p^k$, all such subgroups are conjugate, and their number $n_p$ satisfies
\[
 n_p\equiv1\pmod p, \qquad n_p\mid m.
\]
\end{theorem}

\begin{theorem}[Fundamental Theorem of Finite Abelian Groups]
Every finite abelian group is isomorphic to a direct product of cyclic groups of prime-power order. Equivalently, it is isomorphic to
\[
 C_{d_1}\times\cdots\times C_{d_r},
 \qquad d_1\mid d_2\mid\cdots\mid d_r,
\]
with the invariant factors $d_i>1$ uniquely determined (the trivial group corresponds to an empty product).
\end{theorem}


\begin{theorem}[Sign and Cycle Decomposition]
Every permutation $\sigma\in S_n$ is a product of disjoint cycles, uniquely up to their ordering and cyclic rotation. If its disjoint-cycle decomposition has $c$ cycles, including fixed points, then
\[
 \operatorname{sgn}(\sigma)=(-1)^{n-c}.
\]
\end{theorem}

\begin{definition}[Ring, Integral Domain, and Field]
Unless stated otherwise, rings are associative with an identity and need not be commutative. An integral domain is a nonzero commutative ring with identity and without zero divisors. A field is a nonzero commutative ring in which every nonzero element is invertible.
\end{definition}

\begin{definition}[Ideal, Quotient Ring, and Characteristic]
An ideal $I$ of a ring $R$ means a two-sided ideal. The quotient $R/I$ is a ring. The characteristic $\charac R$ is the additive order of $1_R$, or $0$ if it has infinite additive order. Ring homomorphisms preserve identities.
\end{definition}

\begin{theorem}[Finite Integral Domains]
Every finite integral domain is a field. A finite field has cardinality $p^k$ for some prime $p$ and $k\ge1$, and for each such prime power $q$ there is, up to isomorphism, a unique field $\mathbb F_q$.
\end{theorem}

\begin{theorem}[Lagrange's Theorem on Finite Subgroups of a Field]
Every finite subgroup $G$ of the multiplicative group $F^\times$ of a field $F$ is cyclic. In particular, the multiplicative group of each finite field is cyclic.
\end{theorem}

\begin{theorem}[Multiplicative Group of a Finite Field]
The group $\mathbb F_q^\times$ is cyclic of order $q-1$. In particular, each $a\in\mathbb F_q$ satisfies $a^q=a$, and for $m\in\N$ the equation $x^m=1$ in $\mathbb F_q$ has exactly $\gcd(m, q-1)$ solutions.
\end{theorem}

\begin{proposition}[Irreducible Polynomial Quotients]
For a field $F$ and nonconstant $f\in F[x]$, the quotient ring $F[x]/(f)$ is a field if and only if $f$ is irreducible over $F$. If $F=\mathbb F_q$ and $\deg f=d$, the quotient has $q^d$ elements.
\end{proposition}

\begin{theorem}[Polynomial Method over Finite Fields]
For a nonzero polynomial $f\in\mathbb F_q[x]$ of degree $d$, the equation $f(x)=0$ has at most $d$ solutions in $\mathbb F_q$. A polynomial function on $\mathbb F_q$ has a unique representative of degree at most $q-1$ modulo $x^q-x$.
\end{theorem}
